\chapter{层次和列项}
\section{有标题的条}
条编号和条标题单独占一行，上下各空半行。条文按段编排，段首空两个汉字，回行时顶格。

\subsection{条的细分}
需要进一步细分时可使用带编号的子条，但不应因此增加理解文件的难度。

\subsubsection{子条的标题}
子条的编号、标题及字体与条保持一致。

\stdclause{这是无标题条的条文。条编号之后空一个汉字，条文直接接排。}

\section{列项}
\subsection{破折号列项}
第一层次列项及第二层次列项的缩进和回行位置如下：\index{列项}
\begin{stdlist}
    \item 第一层次列项甲；
    \item 第一层次列项乙：
    \begin{stdlist}
        \item 第二层次列项甲；
        \item 第二层次列项乙，内容较长时用于检查回行后的悬挂缩进位置。
    \end{stdlist}
    \item 第一层次列项丙。
\end{stdlist}

\subsection{字母和数字编号列项}
需要识别或表明先后顺序时，可使用编号列项：
\begin{stdenum}
    \item 编号列项甲；
    \item 编号列项乙：
    \begin{stdenum}
        \item 第二层次数字编号甲；
        \item 第二层次数字编号乙。
    \end{stdenum}
    \item 编号列项丙。
\end{stdenum}

\chapter{符号和缩略语}
下列符号和缩略语适用于本文件。\index{符号和缩略语}
\begin{stdlist}
    \item[$E$] 能量；
    \item[$m$] 质量；
    \item[GB] 国家标准代号；
    \item[ICS] 国际标准分类号。
\end{stdlist}

\chapter{分类和编码}
本章为分类和编码要素的占位示例。正式文件应给出标准化对象的划分以及分类结果的命名或编码规则。\index{分类和编码}

\chapter{总体原则和总体要求}
本章为总体原则和总体要求要素的占位示例。总体原则宜使用陈述型或推荐型条款，总体要求应使用要求型条款。\index{总体要求}

\chapter{核心技术要素}
本章为核心技术要素的占位示例。正式文件应根据标准的功能类型选择需要规定的技术内容。\index{核心技术要素}

\chapter{其他技术要素}
本章为其他技术要素的占位示例。需要设置试验方法、检验规则或标志、包装、运输和贮存等内容时，可在相应章节中编写。\index{其他技术要素}
